% !TeX root = ../papers/nc-p3.tex
\section{混合窓を通る判別層と十五次元成分への特殊化}
\label{sec:marked-mixed-component}
\begin{theorem}\label{thm:marked-mixed-component}
$G,F$ が対称側、$Q,P$ が交代側である有限の混合窓は、印を忘れても、
定理\ref{thm:full-fifteen-component}の十五次元成分に全て属する。
この向きの混合点を通る片側対称な印付き表示では、$g=5,4,3$ の同型類の像を
含む入れ子の既約閉包があり、その次元はそれぞれ $14,12,8$ である。
これらの一般点は全次数で正則で、最小周期は順に $4,4,2$ である。
定理\ref{thm:marked-prescribed-discriminant-as}の指定法線・指定直線の全構成は第二の閉包に属する。
この記述における局所特異点は印と正規化座標を取った表示のものである。
\end{theorem}
\begin{proof}[三層の局所環と閉包]
$S=\mathrm{Sym}_0(V)$ と置き、$G$ は $\bigwedge^2 G'\to\mathfrak{so}(V)$ が可逆な
$S$ の三基底列を動く。ここで $G'=\langle G_i\rangle$ である。
三列 $F$ から $Z\in\mathrm{Hom}(k^3,S/G')$ を取る。交換子の跡条件
$\sum[G_i,F_i]=0$ は、$F_i$ に $G_j$ の倍を加えることで解ける。
補正係数の交代部分に関する微分は三交換子を列に持つ同型である。
よって局所的に十五個の $G$ 座標と任意の $2\times3$ 行列 $Z$、六個の対称な補正係数が自由であり、
$g=3+\operatorname{rank}Z$ となる。$Q=P$ の復元は交代な基点でエタールなので、
完備局所環でもこの積表示が成り立つ。$g\le4$ は滑らかな二十一次元因子と
階数1以下の $2\times3$ 行列錐の積で、特異集合は $Z=0$ である。
実際、階数1なら一つの単元成分で残りを解ける。原点では全小行列式が二次である。
錐は $Z=u v^t$ の座標環で、二次小行列式は同じ行・列多重次数の単項式を
添字交換で全て同一視するため、その関係イデアルを生成する。
この半群は非負で行と列の総次数が等しい全ての次数からなり、飽和しているので錐は正規である。
像直線と核二平面をともに記憶する表示は $\mathbf P^1\times\mathbf P^2$ 上の
$\mathcal O(-1,-1)$ の全空間であり、原点上に零切断を持つ。
像直線からは法線 $N$ が決まり、核二平面の $F$ による像からは共通二平面と
その交換子の直線 $\ell$ が決まる（ランク1で $F$ が独立な開集合）。
従って固定した二因子を動かした前節までの構成もこの同じ局所的な層に入る。

全ての有限混合点を同じ閉包で扱えることを確認する。両基底を合同変換で規格化すれば
$F=-G$ で $Q_i=[Le_i]_\times$ とできる。可逆な $L$ の余因子行列が三交換子から
指定されるため、解は既約な $G$ の開集合の上の次数2のエタール被覆をなす。
この被覆が分離しても二成分を $L\mapsto-L$ が交換する。これは残余群の
$B=-I,A=I,s=1$ で実現され、同型類の像は既約である。
$g\le4$ の局所環は上記の正規な錐により各混合点で幾何学的単枝であり、
制限のない片側対称表示もそこでは滑らかである。
従ってこの混合軌跡の各既約成分に沿って、通過する各層の成分は一つで同じである。
実際、有限個の成分との閉交叉が軌跡を覆い、少なくとも一つが軌跡全体を含む。
有限点条件は開で、有限混合候補は定理\ref{thm:mixed-point-frobenius-algebras}により全て正則である。
このうち分裂二十点例で安定化群は有限であり、印を忘れるファイバーも有限である。
正規化表示の三層の次元 $27,25,21$ から群の次元13を除けば主張の像次元になる。
一般点での正則性は相対bar複体の内部次数ごとの特殊化から得る。
$g=4$ の一般最小周期と射影自己同型の自明性は定理\ref{thm:marked-semisimple-period}より、
配置の番号付けと同じ閉条件を使えば $g=5$ の一般点にも従う。
\end{proof}

\begin{proposition}\label{prop:split-mixed-transverse-curve}
分裂混合点 $G=D$, $F=-D$, $Q=P=K$, $K_i=-[e_i]_\times$ を通り、
\[
             G=D+tK,\quad F=-D+tK,\quad Q=K+tD,\quad P=K-tD
\]
と置く。このアフィン直線上の窓が標準Koszul左右AS正則である必要十分条件は
$(t^2+3)(3t^2+1)\ne0$ である。反対階数は $t=0$ で10、$t^2=-1$ で14、他で16となる。
正則点の周期は2を割る。
\end{proposition}
\begin{proof}
標準印付きの交換子式は全 $t$ で成り立つ。式が高々二次なので $t=0,1,-1$ の
三つで係数ごとに確かめられる。星印の六次単位小行列により二切断は階数6である。
$F_i=-G_i^t$, $P_i=-Q_i^t$ なので二位置移動は各補空間の符号を反転すれば四テンソルを保つ。
反対切断を四元群ラベルの総和ごとに四ブロックにするとその行列式は順に
$-12(1+t^2)^2,16t^2,16t^2,16t^2$ である。零色ブロックの星行と星列には
非対角に1が並び、残る三次正方部の非対角成分は $-2(1+t^2)$、対角は零なので $t^2=-1$ で階数2となる。

$c^2=-2(1+t^2)\ne0$ とする。四位置の星座標を全て $c$ 倍し、さらに四位置で順に
\[
 (i,1,i,1),\quad (i,i,1,1),\quad (-i,1,-i,1),\quad (-i,-i,1,1)
 \qquad (i^2=-1)
\]
の対角倍率を掛けると Walsh 表示になる。すなわち $d=b+e$ とすると
成分 $w_{j,j+b,j+d,j+e}$ は $(-1)^{j_1d_0} m_{be}$ である。
欠色は両方とも0であり、残り三色での符号付き平方行列を $c^2$ で割ると
第一行が $(-2(1+t^2),(1-t)^2,(1+t)^2)$ である巡回行列となる。
行和・列和は零で、全二次余因子は $(t^2+3)(3t^2+1)$ に等しい。
欠色のラベルを移す操作はこの階数を保つ。定理\ref{thm:walsh-normal-quadrics}で判定できる。

$t^2=-1$ では、補文字だけからなる全成分が消え、非零成分は星を一つか二つ持つ。
星を $z$ 倍してテンソル全体を $z$ で割った族において $z=0$ には一つ星のみが残る。
上記の位相倍率を掛けると、これは Walsh 表示で $m_{bb}=0$、
非対角係数は上の $m_{be}/c$ の特殊値である。実際 $b\ne e$ の項は
四ラベルを一度ずつ含み、星倍率が正確に一つである。符号付き平方行列の余因子は
$t^2=-1$ で $-4$ となるのでこの特殊窓は全次数正則である。$z\ne0$ の族は
元の窓と同型であり、標準Frobenius条件および内部次数別の完全性を特殊点から
幾何学的一般点へ移して、元の窓での全次数正則性を得る。
\end{proof}

\begin{proof}[十五次元成分の同定]
上記曲線の $c\ne0$ の全点は定理\ref{thm:walsh-full-component}により既知の
十五次元成分の基底変換による像に属する。分裂混合点で他の成分との交叉の可能性を除くため、
命題\ref{prop:twenty-conormal}を適用する。第一辺の二十点を前節の $(a,b)$ とすれば、
第二切断 $N$ の左右の部分評価は階数3で、核点はそれぞれ $(c,b)$ である
（ここで $c$ は位置2の点）。従って二十個の余法線は
$(c\otimes b)\otimes(a\otimes b)$ である。分裂素数1009でその二十行の階数は20、
二切断の各十次元の左右核の積から得る合計二百行のSchur微分の階数は180となる。
前者による二十個の恒等式があるため後者は標数零で180を超えない。
よって分裂混合点の全テンソル軌跡は滑らかな次元76で、通過成分は既知のもの一つである。
混合点を通る既約な $g=5$ の片側対称族の閉包もこれに入り、三層に共通する成分の同定が従う。
さらに反対行列式は曲線上 $t=0$ に正確に六位の零を持つ。核と余核の六次元間に誘導される
一次写像が可逆であり、この曲線は片側対称な階数十の層から横断方向へ出る。
\path{computations/verify_mixed_component_smoothing.py}が二つの微分、余法線、三値の
交換子、四ブロックとWalshの全係数恒等式を再現する。
\end{proof}
